%% ch04.tex — Chapter 4: The MapReduce Paradigm and the Hadoop Processing Stack
%% Part of "Data Engineering" textbook

\chapter{The MapReduce Paradigm and the Hadoop Processing Stack}
\label{ch:mapreduce}

\epigraph{\textit{``Our abstraction is inspired by the map and reduce primitives present
in Lisp and many other functional languages. We realized that most of our
computations involved applying a map operation to each logical record in
our input in order to compute a set of intermediate key/value pairs, and
then applying a reduce operation to all the values that shared the same key.''}}{Dean \& Ghemawat, OSDI 2004}

\vspace{4pt}

\begin{chapterlearning}
After completing this chapter, you will be able to:
\begin{enumerate}
  \item Explain why a new computation model was required for petabyte-scale
        data on distributed file systems, and identify the specific
        limitations of SQL and MPI that MapReduce addresses.
  \item Write a MapReduce program by defining the Map and Reduce functions,
        and trace the complete execution through Map $\to$ Shuffle/Sort
        $\to$ Reduce using a concrete example.
  \item Explain the roles of the Combiner (local optimisation) and the
        Partitioner (reducer assignment) and quantify the network savings
        a Combiner provides.
  \item Describe the full MapReduce execution pipeline: input splits,
        task scheduling with data locality, speculative execution, and
        output materialisation.
  \item Explain how MapReduce achieves fault tolerance through
        deterministic task re-execution.
  \item Write Pig Latin scripts using LOAD, FILTER, FOREACH, GROUP,
        JOIN, and STORE, and explain how Pig Latin is compiled to
        MapReduce.
  \item Write HiveQL queries using CREATE TABLE, partitioning,
        bucketing, and SELECT with GROUP BY, and explain the role of
        the Hive Metastore.
  \item Describe the HBase data model (row key, column families,
        column qualifiers, timestamps) and the HBase architecture
        (HMaster, RegionServers, WAL, MemStore, HFile).
  \item Select the appropriate tool from the Hadoop stack---raw
        MapReduce, Pig, Hive, or HBase---for a given data engineering
        problem.
\end{enumerate}
\end{chapterlearning}

\bigskip

Chapter~\ref{ch:hdfs} established how petabytes of data are stored
reliably across a cluster of commodity servers. Storage, however, is
only half the challenge. The data must be processed: indexed, aggregated,
joined, filtered, and analysed. In 2003, Google's engineers faced
precisely this problem. Their GFS cluster held hundreds of terabytes of
web crawl data; building the search index required processing every
document, aggregating term frequencies, and computing PageRank---
computations that would take thousands of hours on a single machine.

The solution they published in 2004 is the subject of this chapter:
\term{MapReduce}, a programming model that allows a programmer to express
a large-scale computation as two simple functions---\emph{map} and
\emph{reduce}---and have the framework automatically parallelise the
execution across hundreds or thousands of machines, handling data
partitioning, scheduling, fault tolerance, and inter-machine
communication transparently \parencite{MapReduce2004}. Over the following
decade, MapReduce became the computational backbone of an entire
ecosystem of data processing tools: Apache Pig, Apache Hive, and Apache
HBase, all of which this chapter also covers.

%%====================================================================
\section{Why a New Computation Model?}
\label{sec:mr-motivation}
%%====================================================================

The data processing challenges of petabyte-scale distributed storage
do not yield to the tools that work well for smaller datasets. Two
dominant paradigms---relational SQL and MPI-style parallel
programming---both fail in this environment for distinct reasons.

\textbf{SQL on a single server} requires data to be loaded into a
relational schema, supports random access via indexes, and executes on
a single instance (or a tightly coupled SMP cluster). None of these
properties hold for unstructured petabyte data stored on thousands
of commodity DataNodes. Sequential full-table scans of a petabyte
dataset at 200\MB/s on a single server would take over 57 days.
SQL's assumption of a schema-on-write data model is incompatible with
semi-structured and unstructured data (web pages, log files, binary
media). SQL-on-Hadoop tools such as Hive (discussed in
Section~\ref{sec:hive}) resolve these issues by compiling SQL into
MapReduce jobs.

\textbf{MPI (Message Passing Interface)} is the standard for
scientific high-performance computing. MPI programmes explicitly manage
inter-process communication, data partitioning, and synchronisation.
This gives the programmer maximum flexibility but at enormous cost: the
programmer must manually handle every aspect of distribution. More
critically, MPI does not handle fault tolerance: if any MPI process
fails, the entire job must typically be restarted from scratch. In a
cluster of thousands of commodity machines where failures occur daily,
a long-running MPI job that restarts on every node failure is
effectively unusable.

MapReduce resolves both limitations. It enforces a structured
programming model (Map and Reduce functions) that is restrictive
enough for the framework to automate parallelisation, scheduling, and
fault recovery---but expressive enough to implement a broad class of
large-scale data transformations. The framework, not the programmer,
decides how many tasks to spawn, which machine runs each task, how
to partition intermediate data, and what to do when a machine fails.

\begin{keyinsight}{The Expressiveness of Map + Reduce}
The surprising result of Dean and Ghemawat's paper is that an
enormous variety of large-scale computations---word frequency,
inverted index construction, PageRank, distributed sort, join
operations, machine learning feature extraction---can be expressed
as a single Map phase followed by a single Reduce phase (or a
sequence of MapReduce stages). The restrictive two-function interface
is a \emph{feature}, not a limitation: its uniformity is exactly what
allows the framework to manage all distribution and fault tolerance
automatically.
\end{keyinsight}

%%====================================================================
\section{The MapReduce Programming Model}
\label{sec:mr-model}
%%====================================================================

MapReduce adopts the functional programming concepts of \emph{map} and
\emph{reduce} from Lisp and ML, applied to distributed key-value pairs
at petabyte scale. The programmer provides two pure functions---with no
shared mutable state and no side effects---and the framework handles
everything else.

\begin{defbox}{The MapReduce Functions}
Given input data represented as a set of key-value pairs $(k_1, v_1)$:

\medskip
\noindent\textbf{Map:}
\[
\text{map} : (k_1,\, v_1) \;\longrightarrow\; \text{list}\bigl((k_2,\, v_2)\bigr)
\]
The Map function processes each input record independently and emits
zero or more intermediate key-value pairs.

\medskip
\noindent\textbf{Reduce:}
\[
\text{reduce} : (k_2,\, \text{list}(v_2)) \;\longrightarrow\; \text{list}\bigl((k_3,\, v_3)\bigr)
\]
The Reduce function receives all intermediate values associated with
a single key $k_2$ and aggregates them into the final output.
\end{defbox}

Between the Map and Reduce phases, the framework executes the
\term{shuffle and sort}: it collects all intermediate $(k_2, v_2)$
pairs emitted by all Map tasks, groups them by key, and sorts each
group, delivering to each Reduce task an iterator over all values
$\text{list}(v_2)$ for a single key $k_2$. The programmer never
writes shuffle-and-sort logic; it is the framework's responsibility.

%--------------------------------------------------------------------
\subsection{The Canonical Example: Word Count}
\label{subsec:word-count}
%--------------------------------------------------------------------

The canonical MapReduce example is word count: given a collection of
text documents, compute the frequency of every distinct word across
the entire collection. This example is simple enough to follow in
detail, yet it illustrates every essential phase of MapReduce execution.

\begin{lstlisting}[style=pseudocode, caption={MapReduce word count: Map and Reduce functions}]
# ----- Map function -----
# Input:  (docid, document_text)
# Output: (word, 1) for every word in the document

def map(docid, text):
    for word in text.split():
        emit(word.lower(), 1)

# Example: map("doc1", "To be or not to be") emits:
#   ("to", 1), ("be", 1), ("or", 1), ("not", 1), ("to", 1), ("be", 1)

# ----- Shuffle and Sort (framework-managed) -----
# Groups all emitted pairs by key, sorts alphabetically:
#   ("be",  [1, 1, 1, ...])    <- all 1s from every document
#   ("not", [1, 1, ...])
#   ("or",  [1, ...])
#   ("to",  [1, 1, 1, ...])

# ----- Reduce function -----
# Input:  (word, iterator_of_counts)
# Output: (word, total_count)

def reduce(word, counts):
    total = sum(counts)
    emit(word, total)

# Example: reduce("be", [1,1,1,1,1]) emits: ("be", 5)
\end{lstlisting}

Each document in the collection is assigned to an independent Map
task; if there are 10 million documents, 10 million Map tasks run in
parallel (subject to cluster capacity). Each Map task processes only
its own documents and emits word counts for that subset. The framework
then collects all $(word, 1)$ pairs from all Map tasks, groups them
by word, and routes each group to a single Reduce task, which sums
the counts. The output is the complete word frequency table for the
entire collection.

%--------------------------------------------------------------------
\subsection{Additional MapReduce Examples}
\label{subsec:mr-examples}
%--------------------------------------------------------------------

Word count is illustrative, but the MapReduce model extends to a much
wider class of computations. Two further examples demonstrate its
generality.

\textbf{Inverted index construction.} A search engine's inverted index
maps each term to the list of documents containing it. The Map
function emits $(\text{term}, \text{docid})$ for each term in each
document. After shuffle and sort, each Reduce task receives
$(\text{term}, [\text{docid}_1, \text{docid}_2, \ldots])$ and
concatenates the list into the posting list for that term.

\textbf{Distributed sort.} To sort a petabyte dataset, the Map
function emits each record as $(\text{sort\_key},\, \text{record})$.
The shuffle and sort phase sorts all intermediate keys globally (using
a Partitioner that distributes key ranges uniformly across Reducers).
Each Reducer receives records in sorted order for its key range; the
concatenation of all Reducer outputs is globally sorted. This pattern---
TeraSort---is the standard benchmark for distributed sorting, and
Hadoop clusters have set world records for sorting 1 TB in minutes.

%%====================================================================
\section{The Combiner and the Partitioner}
\label{sec:combiner-partitioner}
%%====================================================================

Two optional but important components sit between the Map and Reduce
phases: the Combiner and the Partitioner. Both are performance
optimisations that the programmer can provide without changing the
correctness of the computation.

%--------------------------------------------------------------------
\subsection{The Combiner: Local Aggregation}
\label{subsec:combiner}
%--------------------------------------------------------------------

In the word count example, each Map task might process a 128\MB\ input
split containing hundreds of thousands of words. If the word ``the''
appears 50,000 times in that split, the Map task emits 50,000
\texttt{("the", 1)} pairs, all of which must be transferred over the
network to the Reducer responsible for ``the''. This is wasteful: the
network cost is proportional to the number of word occurrences, not
the number of distinct words.

The \term{Combiner} is a locally executed mini-Reducer that runs
immediately after the Map phase \emph{on the same node}. It receives
the Map task's output and aggregates values for each key before they
are shuffled across the network. For word count, the Combiner
would reduce the 50,000 \texttt{("the", 1)} pairs to a single
\texttt{("the", 50000)}, reducing the shuffle volume for this key
by a factor of 50,000.

Combiners can be applied whenever the Reduce function is
\emph{commutative} and \emph{associative}: if $f(a, f(b, c)) = f(f(a, b), c)$
and $f(a, b) = f(b, a)$, then partial application of $f$ is always
safe. Sum, count, min, and max all satisfy this property. Average does
not---the Combiner cannot compute a local average because it does not
know how many records the Reducer will receive from other Map tasks.

In practice, enabling a Combiner can reduce shuffle network traffic
by 80--95\% for aggregation workloads, dramatically reducing job
completion time on bandwidth-constrained clusters.

%--------------------------------------------------------------------
\subsection{The Partitioner: Directing Keys to Reducers}
\label{subsec:partitioner}
%--------------------------------------------------------------------

The \term{Partitioner} determines which Reducer receives each
intermediate key. Given a key $k$ and the number of Reduce tasks $R$,
the default HashPartitioner computes:
\begin{equation}
\text{reducer\_id}(k) = h(k) \bmod R
\label{eq:hash-partitioner}
\end{equation}
where $h(k)$ is the hash of key $k$. This distributes keys uniformly
across all Reduce tasks in expectation---the fundamental requirement
for preventing \term{reducer skew} (one reducer receiving far more
data than others and becoming the job's bottleneck).

For workloads where the hash distribution is insufficient---for
example, when sorting requires that all keys in a range go to a
specific Reducer---the programmer provides a custom Partitioner.
The TeraSort example uses a \term{TotalOrderPartitioner} that
divides the keyspace into $R$ equal-sized ranges based on a sample
of the input, ensuring that each Reducer receives a balanced
partition of the sorted key space.

%%====================================================================
\section{The Full MapReduce Execution Pipeline}
\label{sec:mr-execution}
%%====================================================================

Having defined the programming model, we can now trace the complete
execution of a MapReduce job from input files on HDFS to output files
on HDFS, through the framework's scheduling, data-locality optimisation,
and fault recovery mechanisms.

%--------------------------------------------------------------------
\subsection{Input Splits and the InputFormat}
\label{subsec:input-splits}
%--------------------------------------------------------------------

A MapReduce job begins with the client submitting a job to the cluster.
The framework reads the input files from HDFS and divides them into
\term{input splits}---logical chunks of input data, typically aligned
to HDFS block boundaries (128\MB\ by default). One Map task is
instantiated per input split; the number of Map tasks is therefore
approximately $\lceil \text{total input size} / \text{block size}
\rceil$, matching the number of HDFS blocks that contain the input.

The \term{InputFormat} is the abstraction that translates raw HDFS
bytes into $(k_1, v_1)$ pairs that the Map function consumes. The
standard \texttt{TextInputFormat} produces \texttt{(byte\_offset, line)}
pairs---one per line of text. \texttt{SequenceFileInputFormat} reads
Hadoop SequenceFile format. \texttt{CombineFileInputFormat} merges
many small files into a single split, mitigating the small file
problem in the context of MapReduce job scheduling.

%--------------------------------------------------------------------
\subsection{Task Scheduling and Data Locality}
\label{subsec:scheduling}
%--------------------------------------------------------------------

The cluster \term{JobTracker} (in Hadoop 1.x) or YARN
\term{ResourceManager} (Hadoop 2.x and later, covered in
Chapter~\ref{ch:yarn}) accepts the submitted job and begins scheduling
Map tasks. For each input split, the framework attempts to schedule the
Map task on one of the three DataNodes that hold a replica of the
corresponding HDFS block---applying the data locality principle
introduced in Chapter~\ref{ch:foundations}. If a data-local slot is
unavailable, the framework falls back to rack-local, then to
off-rack scheduling.

Data locality is critical for MapReduce performance. A Map task reading
its input from a local disk can sustain 100--300\MB/s of throughput;
the same task reading from a remote rack over a 1 Gbit/s network
connection is limited to approximately 125\MB/s, and in a congested
cluster may receive less. For a job with thousands of Map tasks all
reading simultaneously, the difference between local and off-rack
I/O can mean the difference between a one-hour job and a three-hour
job.

Reduce tasks have no concept of data locality with respect to the
original input---their input comes from the network (the shuffle
phase). Reduce tasks are scheduled on any available node with
sufficient memory and CPU capacity.

%--------------------------------------------------------------------
\subsection{The Shuffle and Sort in Detail}
\label{subsec:shuffle-sort}
%--------------------------------------------------------------------

The shuffle is the most network-intensive phase of MapReduce. As each
Map task produces output, it does not write it directly to HDFS
(which would incur expensive replication); instead, it writes
intermediate data to a local \term{spill buffer} in RAM (default
100\MB). When the buffer is 80\% full, a background thread
\term{spills} the buffer content to disk after sorting by key within
the spill and applying the Combiner (if present). Multiple spills
from a single Map task are merged into a single sorted file before the
task completes.

Each Reduce task then actively \term{fetches} its portion of the
intermediate data from each Map task's output file via HTTP (a process
called the \term{copy phase}). As data arrives from different Map
tasks, the Reduce side merges the multiple sorted streams (using a
merge-sort) to produce a single sorted sequence of $(k_2, v_2)$ pairs
for each key. This merge-sort is the ``sort'' in shuffle-and-sort, and
it is what guarantees that the Reduce function sees values for a single
key in a well-defined, contiguous iterator.

\begin{figure}[h!t]
\centering
\begin{tikzpicture}[
    phase/.style={rectangle, rounded corners=4pt, minimum width=2.5cm,
                  minimum height=1.0cm, font=\small\bfseries\sffamily,
                  align=center},
    arr/.style={-{Stealth[length=5pt]}, very thick},
    lbl/.style={font=\scriptsize\sffamily, align=center}
  ]
  % Input
  \node[phase, fill=DEGray!20, draw=DEGray!50] (input) at (0, 2.0) {HDFS Input\\Blocks};

  % Map tasks
  \node[phase, fill=DELightBlue,  draw=DEBlue!50]  (m1) at (3.2, 3.5) {Map Task 1};
  \node[phase, fill=DELightBlue,  draw=DEBlue!50]  (m2) at (3.2, 2.0) {Map Task 2};
  \node[phase, fill=DELightBlue,  draw=DEBlue!50]  (m3) at (3.2, 0.5) {Map Task 3};

  % Combiner (optional)
  \node[phase, fill=DELightAccent, draw=DEAccent!60, minimum width=2.0cm,
        minimum height=0.6cm, font=\scriptsize\bfseries\sffamily]
       (comb) at (5.4, 2.0) {Combine\\(local)};

  % Shuffle / Sort
  \node[phase, fill=DELightGreen, draw=DEGreen!60] (shuf) at (7.8, 2.0)
        {Shuffle\\+ Sort};

  % Reduce tasks
  \node[phase, fill=DELightPurple, draw=DEPurple!50] (r1) at (10.4, 3.0) {Reduce\\Task A};
  \node[phase, fill=DELightPurple, draw=DEPurple!50] (r2) at (10.4, 1.0) {Reduce\\Task B};

  % Output
  \node[phase, fill=DEGray!20, draw=DEGray!50] (out) at (13.0, 2.0) {HDFS Output};

  % Arrows
  \draw[arr, DEGray!60] (input) -- (m1);
  \draw[arr, DEGray!60] (input) -- (m2);
  \draw[arr, DEGray!60] (input) -- (m3);
  \draw[arr, DEBlue!60] (m1) -- (comb);
  \draw[arr, DEBlue!60] (m2) -- (comb);
  \draw[arr, DEBlue!60] (m3) -- (comb);
  \draw[arr, DEAccent!70] (comb) -- (shuf)
        node[midway, above, lbl] {network\\transfer};
  \draw[arr, DEGreen!70] (shuf) -- (r1);
  \draw[arr, DEGreen!70] (shuf) -- (r2);
  \draw[arr, DEPurple!60] (r1) -- (out);
  \draw[arr, DEPurple!60] (r2) -- (out);

  % Phase labels
  \node[lbl, text=DEBlue]   at (3.2, -0.4) {Map phase};
  \node[lbl, text=DEAccent] at (5.4, -0.4) {Combine};
  \node[lbl, text=DEGreen]  at (7.8, -0.4) {Shuffle \& Sort};
  \node[lbl, text=DEPurple] at (10.4,-0.4) {Reduce phase};
\end{tikzpicture}
\caption{The full MapReduce execution pipeline. Map tasks run in
         parallel, each processing one input split. The optional
         Combiner reduces shuffle traffic. The framework-managed
         shuffle groups and sorts all intermediate data by key before
         delivering it to Reduce tasks.}
\label{fig:mr-pipeline}
\end{figure}

%%====================================================================
\section{Fault Tolerance in MapReduce}
\label{sec:mr-fault}
%%====================================================================

MapReduce's fault tolerance model is both simple and elegant: all
Map and Reduce tasks are \term{deterministic} and \term{idempotent}.
A Map task that processes the same input split will always produce
the same output; a Reduce task that processes the same sorted key-value
sequence will always produce the same output. This determinism is the
foundation of fault recovery.

The \term{JobTracker} (or YARN ResourceManager) monitors the progress
of every running task through periodic \term{heartbeats}. If a task
has not reported progress within a configurable timeout (default
10 minutes), it is considered failed. The framework re-schedules the
failed task on a different node, which re-reads the input split from
HDFS and re-executes the Map function from scratch. Because HDFS
provides three replicas of every input block, the task can always
find its input on a surviving DataNode.

For Reduce tasks, re-execution is more expensive because the Reduce
task may have already fetched a significant portion of its shuffle
data. When a Reduce task fails, the framework restarts it from the
beginning of the copy phase; all already-fetched intermediate data
is discarded and re-fetched from Map task outputs.

\term{Speculative execution} addresses a different failure mode: the
slow task (\term{straggler}). In a large cluster, statistical
variation in hardware performance means that some tasks will always
run significantly slower than average, even without a hard failure.
A single slow Map task on a degraded disk or a hot machine can
delay the entire job, because the Reduce phase cannot begin until
all Map tasks complete. MapReduce mitigates this by detecting tasks
that are running slower than expected and launching \term{backup
tasks}---duplicate executions of the same task on different nodes.
Whichever task---original or backup---completes first writes its
output; the other is killed. This has no correctness impact because
both produce identical outputs (determinism), but can reduce job
completion time by 40\% in practice.

\begin{table}[h!t]
\centering
\caption{MapReduce fault tolerance mechanisms}
\label{tab:mr-fault}
\renewcommand{\arraystretch}{1.4}
\begin{tabular}{L{3.0cm} L{3.4cm} L{4.0cm}}
\toprule
\textbf{Failure type} & \textbf{Detection} & \textbf{Recovery mechanism} \\
\midrule
Map task failure   & Heartbeat timeout       & Re-schedule on another node; re-read HDFS block \\
Reduce task failure & Heartbeat timeout       & Re-schedule; re-fetch all shuffle data \\
TaskTracker failure & Heartbeat timeout       & Re-schedule all tasks that were running on that node \\
Slow task (straggler) & Progress rate below mean & Launch speculative backup task; keep first to complete \\
JobTracker failure  & Not auto-recovered (Hadoop 1.x) & Manual restart; job history lost unless YARN HA used \\
\bottomrule
\end{tabular}
\end{table}

%%====================================================================
\section{Apache Pig: Dataflow Scripting}
\label{sec:pig}
%%====================================================================

Raw MapReduce is powerful but verbose. A simple join of two datasets
requires writing multiple Java classes: an InputFormat for each source,
a custom Mapper that tags records with their source, a custom Partitioner,
and a Reducer that handles the merge logic---hundreds of lines of
boilerplate code. Apache Pig addresses this by providing a high-level
dataflow scripting language called \term{Pig Latin}, which compiles
to one or more MapReduce jobs automatically.

\begin{defbox}{Apache Pig}
Apache Pig is a dataflow platform for analysing large datasets on
Hadoop. Pig Latin programs describe a \textbf{directed acyclic graph
(DAG) of data transformations}: each statement loads, filters,
transforms, groups, joins, or stores a relation. The Pig compiler
translates the DAG into a sequence of MapReduce jobs, optimising
the plan by combining compatible operations into single jobs where
possible.
\end{defbox}

A Pig Latin script operates on \term{relations}---bag-like collections
of tuples, similar to relational tables but without requiring a
predefined schema. The key operators are:

\begin{itemize}
  \item \texttt{LOAD}: read data from HDFS into a named relation.
  \item \texttt{FILTER}: select tuples matching a predicate.
  \item \texttt{FOREACH ... GENERATE}: project or transform fields
        (analogous to SQL's \texttt{SELECT}).
  \item \texttt{GROUP}: group tuples by a key (generates a relation
        of (key, bag) pairs).
  \item \texttt{JOIN}: equi-join two relations on a key.
  \item \texttt{ORDER BY}: sort a relation by a key.
  \item \texttt{STORE}: write a relation to HDFS.
\end{itemize}

The following Pig Latin script analyses a web server access log to
find the most-requested URLs that returned HTTP 404 errors:

\begin{lstlisting}[style=pseudocode, caption={Pig Latin: 404 error frequency analysis}]
-- Load the access log (tab-delimited)
logs = LOAD 'hdfs:///data/access_log'
       USING PigStorage('\t')
       AS (host:chararray, ts:chararray, method:chararray,
           url:chararray, status:int, bytes:long);

-- Keep only 404 responses
errors = FILTER logs BY status == 404;

-- Group by URL
grouped = GROUP errors BY url;

-- Count 404s per URL
counts = FOREACH grouped GENERATE
             group        AS url,
             COUNT(errors) AS num_errors;

-- Sort descending
ranked = ORDER counts BY num_errors DESC;

-- Keep top 20
top20 = LIMIT ranked 20;

-- Write output to HDFS
STORE top20 INTO 'hdfs:///output/top404' USING PigStorage(',');
\end{lstlisting}

The Pig compiler translates this seven-statement script into
typically two MapReduce jobs: one for the FILTER + GROUP + FOREACH
(a shuffle-intensive aggregate), and one for the ORDER BY + LIMIT
(a second sort phase). The programmer writes dataflow logic; the
framework decides how many MapReduce stages are needed and how to
parallelise each one.

Pig is particularly well-suited for ad hoc exploratory analysis on
HDFS data, for ETL pipelines that require complex multi-step
transformations, and for data scientists who are comfortable with
a scripting language but do not want to write Java MapReduce code.
Its primary limitation---shared with raw MapReduce---is high latency:
even simple queries require at least one MapReduce job, with startup
overhead of 30--60 seconds before computation begins. For interactive
analytical queries, Hive with a modern execution engine (Tez or Spark)
or Presto (Chapter~\ref{ch:sql-scale}) is more appropriate.

%%====================================================================
\section{Apache Hive: SQL-on-Hadoop}
\label{sec:hive}
%%====================================================================

Apache Hive, developed at Facebook and open-sourced in 2008, takes a
different approach to Hadoop usability: instead of a new scripting
language, it provides a SQL-like query language called \term{HiveQL}
(or simply ``Hive SQL'') that compiles to MapReduce (and, since Hive
0.13, to Apache Tez or Apache Spark). Hive allows data analysts
comfortable with SQL to query petabyte-scale datasets on HDFS without
learning MapReduce or Pig Latin.

%--------------------------------------------------------------------
\subsection{The Hive Metastore}
\label{subsec:metastore}
%--------------------------------------------------------------------

For SQL queries to work over HDFS files, Hive needs a way to map
SQL table names and column names to HDFS paths and file formats. This
mapping is managed by the \term{Hive Metastore}: a relational database
(typically MySQL or PostgreSQL) that stores table definitions, schema
information, partition metadata, and storage descriptors.

When a user creates a Hive table:

\begin{lstlisting}[style=hiveql, caption={HiveQL: creating a partitioned ORC table}]
CREATE TABLE web_events (
    user_id     BIGINT,
    event_type  STRING,
    url         STRING,
    session_id  STRING,
    duration_ms INT
)
PARTITIONED BY (dt STRING)
STORED AS ORC
LOCATION 'hdfs:///data/web_events';
\end{lstlisting}

Hive records in the Metastore that the table \texttt{web\_events}
maps to the HDFS directory \texttt{/data/web\_events}, is stored
in ORC format, and is partitioned by the column \texttt{dt}
(date string). Data for 2024-01-15 lives in the subdirectory
\texttt{/data/web\_events/dt=2024-01-15/}. The Metastore is the
source of truth that the HiveQL compiler consults to generate the
correct HDFS path for each query.

%--------------------------------------------------------------------
\subsection{Partitioning and Bucketing}
\label{subsec:partitioning-bucketing}
%--------------------------------------------------------------------

\term{Partitioning} is one of the most important performance techniques
in Hive. A partitioned table stores its data in separate HDFS
subdirectories, one per partition value. A query that filters on the
partition column can skip reading all irrelevant partitions entirely---
a technique called \term{partition pruning}. For a table with 365
date partitions and 1\TB\ per partition (365\TB\ total), a query
filtering on a single day reads only 1\TB\ rather than 365\TB:

\begin{lstlisting}[style=hiveql, caption={HiveQL: partition pruning in a query}]
-- Only reads /data/web_events/dt=2024-01-15/
-- Hive's query planner prunes all other 364 partitions
SELECT  event_type,
        COUNT(*)        AS event_count,
        AVG(duration_ms) AS avg_duration_ms
FROM    web_events
WHERE   dt = '2024-01-15'         -- partition filter
  AND   event_type IN ('purchase', 'add_to_cart')
GROUP BY event_type
ORDER BY event_count DESC;
\end{lstlisting}

\term{Bucketing} provides a second level of data organisation within
a partition. Rows are assigned to a fixed number of \term{buckets}
(files) based on the hash of a designated column. Bucketing improves
the efficiency of equi-joins (two tables bucketed on the join key and
with the same number of buckets can be joined without a full shuffle)
and enables efficient sampling (read only one bucket for a 1/$N$
sample).

%--------------------------------------------------------------------
\subsection{HiveQL to MapReduce Compilation}
\label{subsec:hiveql-compilation}
%--------------------------------------------------------------------

When a user submits a HiveQL query, the Hive driver executes the
following steps:
\begin{enumerate}
  \item \textbf{Parse}: convert the SQL string into an abstract syntax
        tree (AST).
  \item \textbf{Semantic analysis}: resolve table and column names
        against the Metastore; validate types and permissions.
  \item \textbf{Logical plan}: convert the AST into a logical query
        plan (a DAG of relational operators: scan, filter, project,
        aggregate, join).
  \item \textbf{Optimisation}: apply rule-based optimisations
        (predicate push-down, column pruning, partition pruning) to
        reduce data read volume.
  \item \textbf{Physical plan}: translate the logical plan into a
        sequence of MapReduce (or Tez, or Spark) jobs.
  \item \textbf{Execution}: submit the jobs to the cluster and monitor
        their progress.
\end{enumerate}

A \texttt{SELECT ... GROUP BY} query maps straightforwardly to a
single MapReduce job: the Map function applies the predicate and
projects the required columns; the Reduce function aggregates by
the GROUP BY key. A multi-table \texttt{JOIN} may generate two or
more MapReduce jobs depending on the join strategy (common-join,
map-join, or bucket-map-join).

\begin{researchspotlight}{Thusoo et al. (2009) --- Hive: A Warehousing Solution over a Map-Reduce Framework}
\textbf{Citation:} Thusoo, A., et~al.\ (2009). Hive -- A Warehousing
Solution over a Map-Reduce Framework. \textit{Proceedings of the
VLDB Endowment}, 2(2):1626--1629. \parencite{Hive2009}

\smallskip\noindent
\textbf{Key insight:} The paper's central argument is that SQL is the
right interface for large-scale data analysis---not because SQL is
technically superior to MapReduce, but because SQL is what hundreds
of millions of analysts already know. Hive's contribution is a
translation layer, not a new programming model.

\smallskip\noindent
\textbf{Technical contributions:}
\begin{itemize}
  \item Defined the Hive data model: tables, partitions, buckets,
        and complex data types (arrays, maps, structs) as SQL
        extensions compatible with semi-structured data.
  \item Introduced the \emph{Metastore} architecture for mapping
        SQL schemas to HDFS directory structures, enabling schema-on-read
        while providing schema-on-write semantics for the SQL layer.
  \item Demonstrated that a HiveQL-to-MapReduce compiler could express
        all SQL-92 analytics workloads (aggregation, joins,
        subqueries, user-defined functions) on petabyte datasets.
\end{itemize}

\smallskip\noindent
\textbf{Impact today:} Hive is the most-used large-scale SQL
engine in the history of open-source data engineering. Its Metastore
is now an independent component used by Spark SQL, Presto/Trino,
and Amazon Athena as their shared table catalog. The HiveQL dialect
has become the \emph{de facto} standard for analytical SQL over
big data platforms, influencing the syntax of SparkSQL, BigQuery,
and Snowflake.
\end{researchspotlight}

%%====================================================================
\section{Apache HBase: Column-Family NoSQL on Hadoop}
\label{sec:hbase}
%%====================================================================

MapReduce, Pig, and Hive are all batch processing tools: they read
from HDFS, process data, and write back to HDFS, with latencies
measured in minutes to hours. There is a class of workload for which
this latency is unacceptable: applications that need to read or write
individual records by key in milliseconds---user profile lookups,
messaging systems, real-time recommendation updates, and fraud
scoring.

Conventional HDFS provides no such capability: its read path requires
a NameNode round trip and a DataNode TCP connection per block,
delivering 50--200 ms latency per operation. Apache HBase, modelled
directly on Google's Bigtable \parencite{Bigtable2006}, provides
\emph{random read/write access} at sub-10 ms latency on data stored
on HDFS---bridging the gap between the Hadoop ecosystem and the
performance requirements of operational applications.

%--------------------------------------------------------------------
\subsection{The HBase Data Model}
\label{subsec:hbase-model}
%--------------------------------------------------------------------

HBase stores data in a model that differs fundamentally from both the
relational model and the key-value model.

\begin{defbox}{The HBase Data Model}
An HBase table is a sorted map of:
\[
(\text{row\_key},\; \text{column\_family}:\text{qualifier},\; \text{timestamp}) \;\longrightarrow\; \text{value}
\]
\begin{itemize}
  \item \textbf{Row key:} an arbitrary byte array, lexicographically
        sorted. Row key design is the central HBase schema design
        decision; the entire table is physically sorted by row key,
        enabling efficient range scans.
  \item \textbf{Column family:} a named group of columns, defined at
        table creation time. All columns in a family are stored
        together on disk.
  \item \textbf{Column qualifier:} a dynamic column name within a
        family, created on the fly without altering the table schema.
        A single row can have zero or thousands of qualifiers within
        a family.
  \item \textbf{Timestamp:} each cell stores multiple versions,
        indexed by timestamp. The most recent version is returned by
        default; older versions are accessible by specifying a
        timestamp.
\end{itemize}
\end{defbox}

Consider a social network's user activity table:

\begin{lstlisting}[style=pseudocode, caption={HBase data model: user activity table}]
# Row structure: (row_key, column_family:qualifier, timestamp) -> value

# User u-10042's profile
("u-10042", "profile:name",     T3) -> "Ada Lovelace"
("u-10042", "profile:country",  T3) -> "GB"
("u-10042", "profile:join_date",T1) -> "2019-03-14"

# User u-10042's activity (sparse: not every qualifier exists for every user)
("u-10042", "activity:last_login",  T5) -> "2024-01-15T14:22:31Z"
("u-10042", "activity:post_count",  T4) -> "847"
("u-10042", "activity:follower_ct", T4) -> "12941"

# Multiple versions: profile:name was updated
("u-10042", "profile:name", T1) -> "Ada Byron"
("u-10042", "profile:name", T3) -> "Ada Lovelace"  <- default (latest)
\end{lstlisting}

The row key design \texttt{u-10042} is intentional: all rows for a
user's profile and activity are co-located in the same row, enabling
an atomic, single-row fetch of the complete user record without a
join. This \term{denormalisation by design} is characteristic of
column-family databases and reflects the NoSQL principle that schema
design should optimise for the read access pattern, not for
normalisation.

%--------------------------------------------------------------------
\subsection{HBase Architecture}
\label{subsec:hbase-arch}
%--------------------------------------------------------------------

HBase's architecture closely mirrors the GFS/HDFS master-worker
pattern, adapted for random access.

The \term{HMaster} is the cluster coordinator. It handles table
creation, deletion, and schema changes; assigns \term{regions}
(contiguous ranges of row keys) to \term{RegionServers}; and
monitors RegionServer health via ZooKeeper. The HMaster is not in
the read/write data path---all client I/O goes directly to
RegionServers after the initial routing lookup.

\term{RegionServers} are the worker processes that handle all
read and write operations. Each RegionServer manages between 10 and
1,000 regions. When a client writes a record, the RegionServer:
\begin{enumerate}
  \item Appends the write to the \term{Write-Ahead Log} (WAL)---an
        append-only HDFS file that ensures durability even if the
        RegionServer crashes before flushing to disk.
  \item Writes the record into the \term{MemStore}: an in-memory
        sorted map (implemented as a Red-Black tree). All reads
        first check the MemStore.
  \item Returns success to the client.
\end{enumerate}

When the MemStore grows beyond a threshold (default 128\MB), it is
\term{flushed} to HDFS as an immutable \term{HFile} (a sorted,
block-compressed file in the HDFS namespace). Reads are served from
a combination of the MemStore (for recent writes) and HFiles (for
older data), with a Bloom filter on each HFile to avoid unnecessary
disk reads for absent keys.

Over time, many small HFiles accumulate for a region. \term{Compaction}
periodically merges HFiles into larger, consolidated files, eliminating
deleted records and expired versions. This is the LSM-tree (Log-Structured
Merge-tree) storage model: by converting all writes into sequential
appends (WAL and HFile flushes), HBase achieves write throughput
competitive with hash-based databases while maintaining sorted order
for efficient range scans.

\begin{researchspotlight}{Chang et al. (2006) --- Bigtable: A Distributed Storage System for Structured Data}
\textbf{Citation:} Chang, F., et~al.\ (2006). Bigtable: A Distributed
Storage System for Structured Data. \textit{Proceedings of the 7th
USENIX OSDI Conference}, pp.\ 205--218. \parencite{Bigtable2006}

\smallskip\noindent
\textbf{Key insight:} Bigtable demonstrates that a single, flexible
data model---the sorted, multi-dimensional map---can efficiently
support workloads as diverse as web indexing (sparse, write-heavy),
Google Earth (large values, read-heavy), and Google Finance
(time-series with many versions). The same model that serves all
of Google's internal applications is now accessible to any developer
through HBase and Google Cloud Bigtable.

\smallskip\noindent
\textbf{Technical contributions:}
\begin{itemize}
  \item Introduced the \emph{column family} abstraction---grouping
        related columns for co-location on disk---which enables
        efficient access to subsets of columns without reading
        entire rows.
  \item Described the \emph{Tablet} architecture (equivalent to
        HBase's Region): range-partitioned, sorted data units
        managed by a Tablet Server, with the master handling
        only assignment and load balancing.
  \item Demonstrated the LSM-tree storage model in a production
        distributed system, establishing the architectural pattern
        that HBase, Cassandra, RocksDB, and LevelDB all follow.
\end{itemize}

\smallskip\noindent
\textbf{Impact today:} Bigtable is cited in over 13,000 papers.
HBase (the open-source implementation) serves as the storage backend
for Apache Phoenix (SQL on HBase), Facebook's internal messaging
infrastructure (1 billion messages per day at peak), and numerous
real-time recommendation systems. Google Cloud Bigtable (the managed
service version) is used by companies including Spotify, Dow Jones,
and T-Mobile.
\end{researchspotlight}

%%====================================================================
\section{Ecosystem Tool Selection}
\label{sec:tool-selection}
%%====================================================================

With four tools in the Hadoop processing stack---raw MapReduce, Pig,
Hive, and HBase---the practical question for a data engineer is which
tool to use for a given problem. The selection matrix below provides
a decision framework based on access pattern, latency requirement,
user skill set, and data structure.

\begin{table}[h!t]
\centering
\caption{Hadoop processing stack tool selection matrix}
\label{tab:tool-selection}
\renewcommand{\arraystretch}{1.55}
\begin{tabular}{L{2.2cm} L{2.6cm} L{1.6cm} L{1.4cm} L{3.0cm}}
\toprule
\textbf{Tool} & \textbf{Best use case} & \textbf{Latency} & \textbf{Interface} & \textbf{Not suited for} \\
\midrule
Raw MapReduce & Custom iterative algorithms; tight performance control & Hours & Java & Interactive queries; rapid development \\
Apache Pig    & Multi-step ETL; ad hoc data exploration; complex transforms & Hours & Pig Latin (script) & Interactive queries; SQL-familiar analysts \\
Apache Hive   & Reporting \& BI; data warehouse queries; SQL analysts & Minutes (batch) & HiveQL (SQL) & Sub-second queries; random reads \\
Apache HBase  & Random read/write; operational data; real-time lookups & $<$10 ms & Java API / REST & Full table scans; complex analytics \\
\bottomrule
\end{tabular}
\end{table}

A practical guideline: use \textbf{Hive} as the default for SQL-capable
analysts running reporting and BI workloads over HDFS data;
use \textbf{Pig} when a pipeline requires complex multi-step ETL
that would be awkward to express in SQL; use \textbf{raw MapReduce}
when performance tuning or an algorithm that does not map cleanly to
Map/Reduce (such as iterative graph algorithms) is required; use
\textbf{HBase} when the application requires sub-millisecond random
access to individual records by key. In practice, Hive over Apache
Tez (which avoids materialising intermediate results to HDFS) or
Hive over Spark (Chapter~\ref{ch:spark}) replaces much of the original
Hive-on-MapReduce workload, delivering 10--100$\times$ lower latency.

%%====================================================================
\section{Case Study: Facebook Data Infrastructure (2008--2012)}
\label{sec:facebook-case}
%%====================================================================

\begin{casestudy}{Facebook --- Hive and HBase at the Petabyte Frontier}

\textbf{Background.}
By 2008, Facebook had accumulated more than 2.5 petabytes of user data
and was generating approximately 15 terabytes of new compressed data
per day. The company needed to provide its data analytics team---several
hundred analysts---with the ability to query this data interactively
in SQL, without requiring them to write MapReduce jobs. Simultaneously,
Facebook Messenger was processing over one billion messages per day and
required sub-millisecond random access to individual message threads
for retrieval.

\smallskip
\textbf{The Hive deployment.}
Facebook's data infrastructure team developed Hive internally and
open-sourced it to Apache in 2008. By 2010, Facebook's Hadoop cluster
had grown to over 2,400 nodes with 30\PB\ of storage, running over
7,500 Hive queries per day submitted by analysts across the company.
The Hive Metastore tracked more than 5,000 registered tables, the
largest of which---Facebook's ``events'' table---exceeded 1\PB\ of
compressed ORC data.

A key engineering challenge was query latency. Early Hive-on-MapReduce
queries over large tables took 20--60 minutes due to MapReduce job
startup overhead and the requirement to materialise all intermediate
results to HDFS. Facebook's team introduced several optimisations
that were later contributed to open source: predicate push-down
(filter early), partition pruning (read only relevant date partitions),
and index-based skip-scan. These reduced the median analyst query
time from 35 minutes to under 5 minutes.

\smallskip
\textbf{The HBase deployment.}
Facebook used HBase as the storage backend for Facebook Messenger
and Inbox. Each user's message thread was stored as an HBase row,
with individual messages as column qualifier-value pairs. The row
key design was critical: Facebook chose a reversed-timestamp prefix
to ensure that the most recent messages for a thread were
lexicographically adjacent on disk, enabling efficient range scans
for loading a conversation inbox.

At peak, Facebook's HBase cluster served over \textbf{one billion
messages per day} across a cluster of several hundred RegionServers.
The P99 read latency for fetching a user's 20 most recent messages
was approximately 8 milliseconds---competitive with purpose-built
NoSQL databases and far below the latency achievable with
HDFS-native batch access.

\smallskip
\textbf{The architectural lesson.}
Facebook's use of both Hive and HBase simultaneously---on the same
underlying HDFS cluster---illustrates a key principle of data
engineering: different workloads require different tools, but those
tools can share the same storage layer. Hive queries and HBase
operations both read from and write to HDFS DataNodes; the
NameNode manages the unified namespace. This shared storage model
avoids data duplication and simplifies the operational burden of
managing two separate storage systems.

By 2012, Facebook's Hadoop cluster had grown to over 100 petabytes,
making it the largest known Hadoop deployment in the world at that
time. The infrastructure team published a series of papers and blog
posts documenting their engineering decisions, which directly
influenced the design of subsequent generations of distributed
analytics tools---particularly the development of Apache Spark, which
addressed the primary limitation of the Hive-on-MapReduce stack:
latency.
\end{casestudy}

%%====================================================================
\section*{Chapter Summary}
\label{sec:ch4-summary}
%%====================================================================

\begin{itemize}
  \item \textbf{MapReduce} resolves two limitations of prior parallel
        computing paradigms: unlike SQL, it requires no predefined
        schema and handles unstructured data on HDFS; unlike MPI, it
        handles fault tolerance automatically through deterministic
        task re-execution.

  \item The MapReduce model requires the programmer to provide two
        functions: \textbf{map} $(k_1, v_1) \to \text{list}(k_2, v_2)$
        and \textbf{reduce} $(k_2, \text{list}(v_2)) \to \text{list}(k_3, v_3)$.
        The framework handles parallelisation, shuffle-and-sort, and
        fault recovery.

  \item The \textbf{Combiner} (local mini-reducer) reduces shuffle
        network traffic for commutative-associative operations. The
        \textbf{Partitioner} controls key-to-reducer assignment via
        hash or custom range logic.

  \item \textbf{Fault tolerance} is achieved through deterministic task
        re-execution on failure; \textbf{speculative execution} launches
        backup tasks for slow stragglers.

  \item \textbf{Apache Pig} compiles a dataflow scripting language
        (Pig Latin) to MapReduce, enabling multi-step ETL without
        Java programming.

  \item \textbf{Apache Hive} compiles HiveQL (SQL dialect) to
        MapReduce. The \textbf{Metastore} maps SQL tables to HDFS paths.
        \textbf{Partitioning} enables partition pruning; \textbf{bucketing}
        accelerates equi-joins and sampling.

  \item \textbf{Apache HBase} provides random read/write access at
        $<$10 ms latency on HDFS, using the LSM-tree model: WAL for
        durability, MemStore for recent writes, HFiles for historical
        data, compaction to merge HFiles.

  \item Tool selection: Hive for SQL analytics; Pig for scripted ETL;
        raw MapReduce for custom algorithms; HBase for random-access
        operational data.
\end{itemize}

%%====================================================================
\section*{Exercises}
\label{sec:ch4-exercises}
%%====================================================================

\exercisesection{Short Questions}

\sq Why does SQL on a single server fail for petabyte-scale data on HDFS?
Name three specific assumptions SQL makes that do not hold in the HDFS
environment.

\sq Write the type signatures of the Map and Reduce functions in the
MapReduce model. What is the framework's responsibility between the Map
and Reduce phases?

\sq Trace the word count MapReduce computation for the input string
``the cat sat on the mat the cat''. List every $(k, v)$ pair emitted
by the Map function, the grouped shuffle output, and the final
$(k, v)$ output of the Reduce function.

\sq What is the Combiner? For what class of Reduce functions can a
Combiner be applied safely? Give an example of a Reduce function for
which a Combiner would give incorrect results.

\sq What is speculative execution in MapReduce? Under what circumstances
does the framework launch a backup task, and why does it not affect
the correctness of the computation?

\sq A MapReduce job has 1,000 Map tasks and 50 Reduce tasks. Describe
what happens when one Map task fails partway through execution.
Describe what happens when one Reduce task fails after it has already
fetched 80\% of its shuffle data.

\sq What is the Hive Metastore? What information does it store, and what
would happen to a HiveQL query if the Metastore database were
unavailable?

\sq Explain partition pruning in Hive. A table has 365 daily partitions
each containing 500\GB\ of ORC data (total 182\TB). A query filters
on a single day. How much data does Hive read, and what is the speedup
factor compared to a full-table scan?

\sq What is the HBase data model? Describe the four coordinates that
uniquely identify a cell in an HBase table.

\sq Explain the LSM-tree write path in HBase: what happens to a client
write at each of the WAL, MemStore, and HFile stages?

\sq What is HBase compaction? Explain the difference between minor
and major compaction, and why compaction is necessary for read
performance.

\sq Compare Apache Pig and Apache Hive on three dimensions: target
user, language paradigm, and primary use case.

\sq A data engineer needs to find the top 10 most active users in a
1-petabyte social media activity dataset (schema: user\_id, event\_type,
ts). Which Hadoop tool would you recommend, and why?

\sq True or false, with justification: ``HBase stores data in HDFS,
so its read latency is determined by HDFS block access latency
(50--200 ms).''\quad\textit{(Hint: consider the MemStore and the
block cache.)}

\bigskip
\exercisesection{Analytical Problems}

\ap \textbf{MapReduce Shuffle Volume Analysis.}
A MapReduce word count job processes a 1\TB\ corpus of English text.
The average word length is 5 characters, and the average word
frequency is 100 occurrences per 128-\MB\ input split. There are
100,000 distinct words in the corpus. The cluster has 500 Map tasks.

\begin{enumerate}[label=(\alph*)]
  \item Estimate the total number of $(word, 1)$ pairs emitted by all
        Map tasks (assume 128\MB\ per split, and that the text has an
        average word density of 150 words per KB).
  \item Estimate the total shuffle data volume in GB without a
        Combiner (each pair requires approximately $|word| + 8$ bytes
        for the key-value pair overhead).
  \item With a Combiner enabled, each Map task emits at most one
        pair per distinct word. Estimate the new shuffle volume and
        the reduction factor.
  \item If the cluster's inter-node network bandwidth is 10 Gbit/s
        per node, and the shuffle phase transfers all data through the
        network simultaneously, estimate the minimum time for the
        shuffle to complete (a) without and (b) with the Combiner.
\end{enumerate}

\ap \textbf{Hive Partitioning Strategy.}
An e-commerce company has a \texttt{transactions} table (500\TB\
total, 5 years of data) with columns: \texttt{transaction\_id},
\texttt{customer\_id}, \texttt{product\_id}, \texttt{amount},
\texttt{country}, \texttt{ts} (timestamp).

\begin{enumerate}[label=(\alph*)]
  \item The company's most frequent query pattern is:
        ``What was the total revenue per product in Germany in Q1 2024?''
        Propose a partitioning scheme. Which column(s) should be the
        partition key(s), and how much data does the optimised query read?
  \item A second query pattern is: ``For a given customer, list all
        their transactions across all time.'' Evaluate whether your
        partitioning scheme from (a) also optimises this query.
        If not, propose how to handle both query patterns (hint:
        consider secondary indexes or a separate HBase table).
  \item If the data is additionally bucketed on \texttt{customer\_id}
        into 256 buckets, and a join query joins this table with a
        30-GB \texttt{customers} dimension table on \texttt{customer\_id},
        explain what bucket-map-join does and estimate the speedup
        over a common join.
\end{enumerate}

\ap \textbf{HBase Row Key Design Analysis.}
A healthcare analytics system stores patient vital signs readings in
HBase. Each reading has: \texttt{patient\_id} (8-char string),
\texttt{metric} (heart\_rate, SpO2, temperature, blood\_pressure),
and \texttt{timestamp} (epoch milliseconds).

The following three row key designs are proposed:

\begin{itemize}
  \item \textbf{Design A:} \texttt{patient\_id + timestamp}
        (e.g., \texttt{P-000042|1705315200000})
  \item \textbf{Design B:} \texttt{timestamp + patient\_id}
        (e.g., \texttt{1705315200000|P-000042})
  \item \textbf{Design C:} \texttt{patient\_id + reverse\_timestamp}
        (e.g., \texttt{P-000042|9994684799999})
\end{itemize}

\begin{enumerate}[label=(\alph*)]
  \item For Design B, describe the write hotspot problem: if all
        writes have approximately the same timestamp, what happens to
        the distribution of writes across RegionServers?
  \item For Design A, explain how a range scan for ``all vitals for
        patient P-000042 in the last 24 hours'' would be executed.
  \item For Design C (reverse timestamp), explain what a scan for
        ``most recent 10 readings for patient P-000042'' looks like.
        Why does this design outperform Design A for this access pattern?
  \item Propose a salting strategy for Design A that eliminates the
        write hotspot for high-frequency patients.
\end{enumerate}

\ap \textbf{MapReduce vs. Pig vs. Hive Performance Comparison.}
A financial services firm must compute the total trading volume per
stock symbol per day for one year of trade records (1.5\TB\, 150
DataNodes, 10 Gbit/s network). Estimate completion time for each tool:

\begin{enumerate}[label=(\alph*)]
  \item \textbf{Raw MapReduce} (Java): Map emits \texttt{(symbol+date, amount)};
        Reduce sums. Assume 80\% data locality, 200\MB/s local read.
        50 Reduce tasks. Estimate total Map I/O time, shuffle time
        (assume 20\% of input volume is shuffled), and Reduce time.
  \item \textbf{Apache Pig}: same logical plan as MapReduce, compiled
        to one MapReduce job with 30 seconds additional startup overhead
        per job. Estimate total elapsed time.
  \item \textbf{Apache Hive on MapReduce}: same as Pig with 60 seconds
        additional startup overhead (Metastore lookup + plan compilation).
        But if the table is partitioned by \texttt{date}, and the query
        selects a 30-day period, how much data is read, and what is the
        new estimated time?
  \item Which tool minimises calendar time? Which minimises engineering
        effort? Which minimises query latency for an analyst who runs
        this query daily?
\end{enumerate}

\bigskip
\exercisesection{Design Problems}

\dpr \textbf{MapReduce Pipeline for a Search Engine Inverted Index.}
Design a MapReduce pipeline that builds a compressed inverted index
from a 500-terabyte web crawl stored in HDFS. Each crawl record is
a JSON document containing the URL, crawl timestamp, and raw HTML.
The inverted index maps each term to a posting list of
\texttt{(url, tf-idf\_score)} pairs, sorted by score descending.

Your design must specify:
\begin{enumerate}[label=(\alph*)]
  \item The number of MapReduce stages required and what each stage
        computes (hint: tf-idf requires knowing the total document
        frequency of each term across the entire corpus, which is not
        available in a single pass).
  \item The Map and Reduce function type signatures and logic for
        each stage.
  \item The Partitioner strategy to ensure all records for the same
        term go to the same Reducer in the appropriate stage.
  \item How the Combiner can be applied in Stage 1 and what the
        network savings are for a typical term with 10,000 occurrences
        per 128-MB input split.
  \item How the pipeline handles the case where a single extremely
        common term (e.g., ``the'') causes a Reducer skew: the single
        Reducer responsible for ``the'' processes 100$\times$ more data
        than the average Reducer.
\end{enumerate}

\dpr \textbf{HBase Schema Design for a Real-Time Recommendation System.}
An online streaming music service needs to store and serve user
listening history and song metadata for a real-time recommendation
engine. Requirements:
\begin{itemize}
  \item 50 million users; each user has an average of 2,000 play
        events, growing at 10 events/day.
  \item Each play event: \texttt{user\_id}, \texttt{song\_id},
        \texttt{play\_ts}, \texttt{duration\_played\_sec},
        \texttt{skipped} (boolean).
  \item The recommendation engine reads the 200 most recent plays
        for a user in $<$5 ms.
  \item A nightly batch job (Hive or Spark) must scan all plays
        in the last 7 days across all users to recompute
        collaborative filtering features.
\end{itemize}

Design the HBase schema. Your design must specify:
\begin{enumerate}[label=(\alph*)]
  \item The table name(s), column families, and column qualifiers.
  \item The row key format, with explicit justification for how it
        supports the 5-ms recommendation read requirement.
  \item How the ``most recent 200 plays'' query is executed
        (scan direction, stop condition, use of reverse timestamp).
  \item How to configure cell versioning to automatically expire
        play events older than 90 days without a separate compaction job.
  \item How the nightly batch Hive/Spark job reads from HBase
        (what API it uses and what the scan looks like).
  \item The estimated RegionServer count if each region is limited
        to 10\GB, and the strategy for pre-splitting regions at
        table creation to avoid write hotspots.
\end{enumerate}

\bigskip
\exercisesection{Programming Exercises}

\pe \textbf{MapReduce Word Count in Python (MRJob-style).}
Implement a word count MapReduce job in pure Python without any
framework, simulating the Map, Shuffle/Sort, and Reduce phases.
Extend it with a Combiner and measure the shuffle volume reduction.

\begin{lstlisting}[style=pyspark, caption={PE4.1: MapReduce word count with Combiner}]
import re
from collections import defaultdict

documents = [
    ("doc1", "To be or not to be that is the question"),
    ("doc2", "All the world is a stage and all the men and women"),
    ("doc3", "To be is to do to do is to be to be is to do"),
    ("doc4", "Be not afraid of greatness some are born great"),
]

# ---------- MAP PHASE ----------
def mapper(doc_id, text):
    """Emit (word, 1) for every word in text."""
    for word in re.findall(r'[a-z]+', text.lower()):
        yield (word, 1)

# ---------- COMBINER (optional local aggregation) ----------
def combiner(map_output):
    """
    Locally aggregate (word, 1) pairs from a single mapper.
    Input:  iterable of (word, 1) tuples
    Output: iterable of (word, local_count) tuples
    """
    local_counts = defaultdict(int)
    for word, count in map_output:
        local_counts[word] += count
    return list(local_counts.items())

# ---------- SHUFFLE AND SORT ----------
def shuffle_and_sort(all_map_output):
    """
    Group all (word, value) pairs by word.
    Returns dict: word -> list of values
    """
    grouped = defaultdict(list)
    for word, val in all_map_output:
        grouped[word].append(val)
    return dict(sorted(grouped.items()))  # sorted by key

# ---------- REDUCE PHASE ----------
def reducer(word, values):
    """Sum all counts for a word."""
    return (word, sum(values))

# ---------- JOB RUNNER ----------
def run_mapreduce(documents, use_combiner=True):
    all_map_output = []
    per_mapper_counts = []

    for doc_id, text in documents:
        map_out = list(mapper(doc_id, text))
        per_mapper_counts.append(len(map_out))

        if use_combiner:
            map_out = combiner(map_out)

        all_map_output.extend(map_out)

    shuffle_out = shuffle_and_sort(all_map_output)

    results = {}
    for word, values in shuffle_out.items():
        _, count = reducer(word, values)
        results[word] = count

    shuffle_volume = sum(len(str(w)) + len(str(v))
                         for w, v in all_map_output)
    return results, shuffle_volume, per_mapper_counts

# Run without and with Combiner
res_no_comb,  vol_no,  cnts_no  = run_mapreduce(documents, use_combiner=False)
res_with_comb, vol_yes, cnts_yes = run_mapreduce(documents, use_combiner=True)

print("Top 10 words:")
for w, c in sorted(res_with_comb.items(), key=lambda x: -x[1])[:10]:
    print(f"  {w:<15} {c}")

print(f"\nShuffle volume WITHOUT combiner: {vol_no:,} bytes")
print(f"Shuffle volume WITH    combiner: {vol_yes:,} bytes")
print(f"Reduction factor: {vol_no / vol_yes:.1f}x")
\end{lstlisting}

\pe \textbf{HiveQL Query Writing and Partitioning Analysis (Python simulation).}
Simulate the effect of Hive partitioning on query I/O by implementing
a partition-pruning model that shows how much data is skipped for
various query predicates.

\begin{lstlisting}[style=pyspark, caption={PE4.2: Hive partition pruning simulator}]
import random
from datetime import date, timedelta

# Simulate a partitioned Hive table: (date, country) -> size_gb
random.seed(42)
countries = ["US", "GB", "DE", "FR", "JP", "IN", "BR", "AU"]
start = date(2023, 1, 1)

# Build partition metadata (what the Hive Metastore stores)
partitions = {}
for i in range(365):
    dt = (start + timedelta(days=i)).isoformat()
    for country in countries:
        size_gb = random.uniform(0.5, 2.5)
        partitions[(dt, country)] = round(size_gb, 3)

total_tb = sum(partitions.values()) / 1024
print(f"Total table size: {total_tb:.2f} TB across "
      f"{len(partitions):,} partitions\n")

def query_io(predicate_dt=None, predicate_country=None):
    """
    Simulate partition pruning for a query with optional filters.
    predicate_dt      : exact date string or None (full scan)
    predicate_country : country code or None (full scan)
    Returns: (partitions_read, gb_read)
    """
    read = {k: v for k, v in partitions.items()
            if (predicate_dt      is None or k[0] == predicate_dt)
            and (predicate_country is None or k[1] == predicate_country)}
    return len(read), sum(read.values())

queries = [
    ("Full table scan",
     dict(predicate_dt=None, predicate_country=None)),
    ("Single date (2024-01-15)",
     dict(predicate_dt="2024-01-15", predicate_country=None)),
    ("Single country (US), all dates",
     dict(predicate_dt=None, predicate_country="US")),
    ("Single date + country (2024-01-15, DE)",
     dict(predicate_dt="2024-01-15", predicate_country="DE")),
]

total_parts = len(partitions)
total_gb    = sum(partitions.values())

print(f"{'Query':<45} {'Parts read':>12} {'GB read':>10} {'Pruned':>8}")
print("-" * 78)
for desc, kwargs in queries:
    n_parts, gb = query_io(**kwargs)
    pruned_pct = (1 - n_parts / total_parts) * 100
    print(f"{desc:<45} {n_parts:>12,} {gb:>10.1f} {pruned_pct:>7.1f}%")
\end{lstlisting}

\pe \textbf{HBase MemStore and Compaction Simulator (Python).}
Simulate HBase's write path: WAL append, MemStore insertion, flush to
HFile, and minor compaction.

\begin{lstlisting}[style=pyspark, caption={PE4.3: HBase write path and compaction simulator}]
import time

MEMSTORE_THRESHOLD = 10  # flush after 10 entries (demo scale)
COMPACTION_THRESHOLD = 3  # compact after 3 HFiles

class HBaseRegionSimulator:
    def __init__(self, region_name):
        self.region_name   = region_name
        self.wal           = []          # Write-Ahead Log (list of ops)
        self.memstore      = {}          # row_key -> {col -> value}
        self.hfiles        = []          # list of sorted dicts
        self.write_count   = 0

    def put(self, row_key, column, value):
        ts = int(time.time() * 1000)
        # 1. Append to WAL
        self.wal.append({"op": "PUT", "row": row_key,
                          "col": column, "val": value, "ts": ts})
        # 2. Write to MemStore
        if row_key not in self.memstore:
            self.memstore[row_key] = {}
        self.memstore[row_key][column] = (value, ts)
        self.write_count += 1

        # 3. Flush if threshold reached
        if len(self.memstore) >= MEMSTORE_THRESHOLD:
            self._flush()

    def _flush(self):
        if not self.memstore:
            return
        # Sort by row key (lexicographic) and materialise as HFile
        hfile = dict(sorted(self.memstore.items()))
        self.hfiles.append(hfile)
        print(f"  [FLUSH] MemStore -> HFile #{len(self.hfiles)} "
              f"({len(hfile)} rows)")
        self.memstore.clear()

        # Compact if too many HFiles
        if len(self.hfiles) >= COMPACTION_THRESHOLD:
            self._minor_compact()

    def _minor_compact(self):
        print(f"  [COMPACT] Merging {len(self.hfiles)} HFiles...")
        merged = {}
        for hfile in self.hfiles:
            for row, cols in hfile.items():
                if row not in merged:
                    merged[row] = {}
                merged[row].update(cols)
        self.hfiles = [dict(sorted(merged.items()))]
        print(f"  [COMPACT] Done. 1 HFile with {len(merged)} rows.")

    def get(self, row_key):
        # Check MemStore first
        if row_key in self.memstore:
            return {"source": "memstore", "data": self.memstore[row_key]}
        # Scan HFiles (most recent first)
        for hfile in reversed(self.hfiles):
            if row_key in hfile:
                return {"source": "hfile", "data": hfile[row_key]}
        return None

    def status(self):
        print(f"\nRegion '{self.region_name}' status:")
        print(f"  WAL entries : {len(self.wal)}")
        print(f"  MemStore    : {len(self.memstore)} rows")
        print(f"  HFiles      : {len(self.hfiles)}")

# Simulate 25 writes to a user-activity HBase region
region = HBaseRegionSimulator("user-activity-0001")

users = [f"u-{i:05d}" for i in range(1, 8)]
import random; random.seed(0)
for _ in range(25):
    uid = random.choice(users)
    col = random.choice(["activity:logins", "activity:purchases",
                          "profile:name", "profile:country"])
    val = str(random.randint(1, 1000))
    region.put(uid, col, val)

region._flush()  # flush remaining MemStore
region.status()

# Test a GET
print("\nGET u-00001:", region.get("u-00001"))
\end{lstlisting}
